\section{Introduction}
\label{sec:intro}

Transformer layers evolve representations through learned additive
updates~\citep{vaswani2017attention,elhage2021mathematical} on top of residual
streams~\citep{he2016deep}. Geometrically, each update performs two distinct
roles relative to the incoming representation: a \textit{parallel component}
that modulates magnitude while preserving direction, and a \textit{perpendicular
component} that redirects it into new semantic subspaces. This dichotomy
reframes representation evolution as a problem of \textit{functional geometry}:
decomposing learned updates by their directional role and testing their
behavioral sensitivity through targeted interventions.

Across diverse pretrained models, our measurements reveal that learned updates
consistently contain substantial components parallel to their incoming states.
Geometrically, parallel updates amount to simple scalar rescaling, which the
residual stream already provides at zero parameter cost~\citep{he2016deep,elhage2021mathematical}.
Yet, heavily parameterized attention and MLP sub-layers actively allocate
capacity to produce such components. This paradox prompts a fundamental question:
are these parallel updates behaviorally essential to Transformer capabilities,
or are they functionally redundant?

To resolve this question, we formulate directional decomposition across two
complementary representation spaces: \textit{residual space}, which analyzes
sub-layer updates relative to the incoming hidden state, and \textit{attention
value space}, which analyzes internal value aggregation relative to the token's
own value representation. Using component-scaling interventions, we systematically
modulate the parallel and perpendicular components in these spaces to measure
their behavioral sensitivity.

Our interventions reveal a pronounced directional asymmetry across both
representation spaces. Perpendicular scaling is consistently disruptive,
confirming that orthogonal steering is essential to model capabilities.
In contrast, parallel scaling is comparatively benign, leaving performance
near baseline across a wide scaling interval. Within this parallel resilience,
the degree of stability depends on the representation space: in attention value
space, scaling the parallel component of cross-token aggregation keeps model
behavior remarkably close to baseline, whereas residual-space parallel scaling is
comparatively less stable. Thus, parallel updates exhibit broad resilience
compared to perpendicular steering, especially within internal value aggregation.

Finally, this directional geometry extends beyond inference-time editing to
post-training compression and from-scratch pretraining. For model
compression~\citep{Frantar2022GPTQAP,Lin2023AWQAW,Sun2023ASA,men2024shortgpt,He2026UncoveringTR},
decomposing update distortion reveals that better-performing quantized and
pruned models consistently exhibit lower perpendicular error, whereas parallel
error is far less discriminative. Motivated by the inference-time asymmetry, we
separately evaluate parallel attention suppression during pretraining. Across
model scales from 300M to 2.7B parameters, the value-space variant consistently
yields lower validation-loss trajectories, and both variants improve average
downstream benchmark scores, with the value-space variant producing the largest
gains.

\begingroup\setlength{\parskip}{0pt}In summary, the contributions of this work are as follows:
\begin{enumerate}
  \setlength{\topsep}{0pt}
  \setlength{\partopsep}{0pt}
  \setlength{\itemsep}{2pt}
  \setlength{\parsep}{0pt}
  \setlength{\parskip}{0pt}
  \item This work reveals substantial direction-preserving components in learned
        transformer updates and maps their behavioral robustness across
        representation spaces.
  \item For attention updates, value-space decomposition separates direct
        self-value flow from cross-token aggregation, revealing a marked
        stability gap between value-space and residual-space edits.
  \item For compression and pretraining, this directional geometry extends beyond
        editing: perpendicular error tracks compressed-model quality, while
        parallel suppression consistently improves validation loss and
        downstream performance across model scales.
\end{enumerate}
\endgroup
